\chapter{二项分布}\label{ch:g11:binom}

独立重复做同一个是/否试验若干次，并统计成功次数：
所得分布 --- 即\emph{二项分布} --- 是
所有离散分布中最重要的一种。本章借助树状图与路径计数构造它；
路径数的闭式公式（含阶乘）
随\cref{ch:g12:comb}中的计数工具一并给出，
并在\cref{ch:g12:randvar}中再次讨论该分布。

\section{伯努利试验}

\begin{definition}[伯努利试验]\label{def:g11:binom:bernoulli}
称恰有两种结果的试验为\emph{伯努利试验}\index{伯努利试验}，其结果为
\emph{成功}（概率为 $p$）与\emph{失败}（概率为 $1 - p$）。
设随机变量 $X$ 在成功时取 $1$，在失败时取 $0$，
则称其服从
\emph{伯努利分布} $\mathcal B(p)$，且
\[
\E(X) = p,
\qquad
\V(X) = p(1 - p).
\]
\end{definition}

\begin{proof}[两个公式的证明]
$\E(X) = p \times 1 + (1-p) \times 0 = p$；又因 $X^2 = X$（$0$
与 $1$ 都与自身的平方相等），故 $\E(X^2) = p$，由
\cref{prop:g11:prob:konig}，$\V(X) = p - p^2 = p(1-p)$。
\end{proof}

\begin{definition}[重复独立试验]\label{def:g11:binom:repeated}
将伯努利试验重复 $n$ 次且\emph{独立}进行，是指：
每次试验的结果不影响其他试验，且任一完整结果序列的概率
等于其在树中对应路径上各概率的\emph{乘积} ---
即每次成功取 $p$，
每次失败取 $1 - p$。
\end{definition}

\begin{example}\label{ex:g11:binom:threetrials}
成功概率为 $p$ 的三次独立试验。序列 SFS
（成功、失败、成功）的概率为 $p(1-p)p = p^2(1-p)$ --- 而
同样，\emph{任一}恰有两次成功的序列，其概率亦为此值，与
成功与失败的位置无关，只有 S 与 F 的\emph{个数}起作用。
\end{example}

\section{路径数与二项式系数}

\begin{definition}[二项式系数]\label{def:g11:binom:coefficient}
在 $n$ 次独立试验的树中，
称\emph{二项式系数}\index{二项式系数} $\binom{n}{k}$（读作“$n$
选 $k$”）为恰含 $k$ 次成功的路径条数。
\end{definition}

\begin{example}
$\binom{3}{2} = 3$：路径为 SSF、SFS、FSS。同理 $\binom{3}{0} = 1$
（路径 FFF），$\binom{3}{1} = 3$ 与 $\binom{3}{3} = 1$。由
约定与树状图可知，$\binom n0 = \binom nn = 1$ 对任意 $n$ 均成立。
\end{example}

\begin{omfigure}
\begin{tikzpicture}[font=\small]
  \coordinate (root) at (0,0);
  \node (S) at (1.5,1.5) {S};
  \node (F) at (1.5,-1.5) {F};
  \node (SS) at (3.0,2.25) {S};
  \node (SF) at (3.0,0.75) {F};
  \node (FS) at (3.0,-0.75) {S};
  \node (FF) at (3.0,-2.25) {F};
  \node (SSS) at (4.5,2.6) {S};
  \node (SSF) at (4.5,1.85) {F};
  \node (SFS) at (4.5,1.1) {S};
  \node (SFF) at (4.5,0.35) {F};
  \node (FSS) at (4.5,-0.35) {S};
  \node (FSF) at (4.5,-1.1) {F};
  \node (FFS) at (4.5,-1.85) {S};
  \node (FFF) at (4.5,-2.6) {F};
  \draw[omThm, thick] (root) -- (S);
  \draw[omThm, thick] (S) -- (SS);   \draw[omThm, thick] (SS) -- (SSF);
  \draw[omThm, thick] (S) -- (SF);   \draw[omThm, thick] (SF) -- (SFS);
  \draw[omThm, thick] (root) -- (F);
  \draw[omThm, thick] (F) -- (FS);   \draw[omThm, thick] (FS) -- (FSS);
  \draw[gray] (SS) -- (SSS);
  \draw[gray] (SF) -- (SFF);
  \draw[gray] (FS) -- (FSF);
  \draw[gray] (F) -- (FF);
  \draw[gray] (FF) -- (FFS);
  \draw[gray] (FF) -- (FFF);
  \node[gray, anchor=west] at (5.0,2.6) {$3$ 次成功};
  \node[omThm, anchor=west] at (5.0,1.85) {$2$ 次成功};
  \node[omThm, anchor=west] at (5.0,1.1) {$2$ 次成功};
  \node[gray, anchor=west] at (5.0,0.35) {$1$ 次成功};
  \node[omThm, anchor=west] at (5.0,-0.35) {$2$ 次成功};
  \node[gray, anchor=west] at (5.0,-1.1) {$1$ 次成功};
  \node[gray, anchor=west] at (5.0,-1.85) {$1$ 次成功};
  \node[gray, anchor=west] at (5.0,-2.6) {$0$ 次成功};
\end{tikzpicture}

{\small $n = 3$ 次试验的树：$\binom{3}{2} = 3$ 条路径（红色）恰含
两次成功，每条路径的概率均为 $p^2(1-p)$。}
\end{omfigure}

\begin{proposition}[帕斯卡法则]\label{prop:g11:binom:pascal}
\index{帕斯卡三角形}
对 $1 \leq k \leq n - 1$，有
\[
\binom{n}{k} = \binom{n-1}{k-1} + \binom{n-1}{k}.
\]
\end{proposition}

\begin{proof}
将 $n$ 次试验的树中恰含 $k$ 次成功的路径，按
其\emph{最后一次}试验的结果分类。以成功结尾的路径，
由前 $n-1$ 次试验中恰含 $k - 1$ 次成功的路径延伸得到，共有
$\binom{n-1}{k-1}$ 条。以失败结尾的路径，由含 $k$ 次成功的
前 $n-1$ 次试验中的路径延伸得到，共 $\binom{n-1}{k}$ 条。
每条路径恰属于两类中的一类。
\end{proof}

帕斯卡法则逐行生成系数 --- 每个数
都是其上方两数之和：
\[
\begin{array}{ccccccccccc}
&&&&&1&&&&&\\
&&&&1&&1&&&&\\
&&&1&&2&&1&&&\\
&&1&&3&&3&&1&&\\
&1&&4&&6&&4&&1&\\
1&&5&&10&&10&&5&&1
\end{array}
\]

\begin{remark}
闭式公式 $\binom nk = \frac{n!}{k!(n-k)!}$ 连同系统的计数理论，
将在\cref{ch:g12:comb}中建立。在本阶段，
帕斯卡三角即可算出所需的每一个系数。
\end{remark}

\section{二项分布}

\begin{theorem}[二项分布]\label{thm:g11:binom:binomial}
设 $X$ 为 $n$ 次独立伯努利试验中成功的次数，
其中每次试验的参数为 $p$，则 $X$ 服从
\emph{二项分布}\index{二项分布} $\mathcal B(n, p)$：
\[
\P(X = k) = \binom{n}{k}\, p^k (1-p)^{n-k},
\qquad k = 0, 1, \dots, n .
\]
\end{theorem}

\begin{proof}
事件 $X = k$ 即恰含 $k$ 次成功的所有路径构成的
每条这样的路径概率为 $p^k(1-p)^{n-k}$：沿路径的乘积
中含 $k$ 个因子 $p$ 与 $n - k$ 个因子 $1-p$，次序任意
（见\cref{def:g11:binom:repeated}）。这样的路径共有 $\binom nk$ 条
（见\cref{def:g11:binom:coefficient}），它们的概率相加即得。
\end{proof}

\begin{example}\label{ex:g11:binom:quiz}
一次测验有 $5$ 道相互独立的题，每题有 $4$ 个选项；一名学生
随机作答，因此每题成功的概率为 $p = \frac14$。答对题数
$X$ 服从 $\mathcal B\left(5, \frac14\right)$，
利用帕斯卡三角第 $5$ 行，有
\[
\P(X = 2) = \binom52 \left(\frac14\right)^2\left(\frac34\right)^3
= 10 \times \frac{1}{16} \times \frac{27}{64}
= \frac{270}{1024} \approx 0.26 .
\]
至少答对一题的概率利用对立事件计算：
$\P(X \geq 1) = 1 - \P(X = 0) = 1 - \left(\frac34\right)^5 \approx
0.76$。
\end{example}

\begin{proposition}[期望与方差]\label{prop:g11:binom:expectation}
若 $X \sim \mathcal B(n, p)$，则：
\[
\E(X) = np,
\qquad
\V(X) = np(1-p).
\]
\end{proposition}

\begin{proof}[说明]
记 $X = X_1 + X_2 + \dots + X_n$，其中 $X_i$ 等于 $1$，
当第 $i$ 次试验成功时；每个 $X_i$ 都是期望为 $p$ 的伯努利变量
（见\cref{def:g11:binom:bernoulli}）。期望可加：把 $n$ 个贡献相加，
得 $\E(X) = np$。对独立变量，\emph{方差}也可相加，这一点
成立，但更为细致：方差公式是
\emph{在本阶段不加证明地承认}，并在\cref{ch:g12:sums}中证明。
\end{proof}

\begin{omfigure}
\begin{tikzpicture}
\begin{axis}[
  width=10cm, height=5cm,
  ybar, bar width=9pt,
  xlabel={$k$}, ylabel={$\P(X = k)$},
  xmin=-0.8, xmax=10.8, ymin=0, ymax=0.3,
  xtick={0,1,2,3,4,5,6,7,8,9,10},
  ytick={0.1, 0.2},
  tick label style={font=\small},
  label style={font=\small},
  title={$\mathcal B(10,\ 0.5)$}]
  \addplot[fill=omDef!40, draw=omDef] coordinates {
    (0,0.00098) (1,0.00977) (2,0.04395) (3,0.11719) (4,0.20508)
    (5,0.24609) (6,0.20508) (7,0.11719) (8,0.04395) (9,0.00977)
    (10,0.00098)};
\end{axis}
\end{tikzpicture}

{\small $\mathcal B(10, 0.5)$ 的分布（十次公平抛硬币）：
中心位于 $\E(X) = np = 5$，图形对称，几乎全部概率
集中在 $2$ 与 $8$ 之间。}
\end{omfigure}

\begin{method}[识别二项情形]\label{met:g11:binom:recognize}
在写出 $X \sim \mathcal B(n, p)$ 之前，需检验三个要素：
\emph{固定的}试验次数 $n$，且事先确定；每次试验有
\emph{两种结果}，且成功概率都是\emph{相同的} $p$；试验是
\emph{独立}的（有放回抽样，或使用各自独立的装置）。
从小总体中不放回抽取，\emph{并非}
二项情形——每次抽取后，概率都会改变
（见\cref{exo:g11:prob:6}）。
\end{method}

\section{抽样：观察结果是否出人意料？}

二项分布回答了一个非常实际的问题：\emph{若成功概率确实为 $p$，
那么成功次数取怎样的值
是合理的？}

\begin{example}\label{ex:g11:binom:decision}
按设计，机器产出的次品至多占 $10\%$。在一
批 $10$ 件产品中有 $4$ 件次品。是运气不好，还是机器坏了？若
机器运转正常，则次品数
服从 $\mathcal B(10, 0.1)$，且有
\[
\P(X \geq 4) = 1 - \P(X \leq 3) \approx 1 - 0.987 = 0.013 :
\]
约为 $80$ 分之一的机会。观察到如此小概率的事件，是一个强烈的
信号 --- 应\emph{拒绝}“机器仍
以 $10\%$ 的次品率运转”这一假设，同时须记住，
该判断出错的概率约为 $0.013$。
\end{example}

\begin{method}[基于二项模型的判定准则]\label{met:g11:binom:decision}
对于观察到的成功次数 $k$，将其与假设
$X \sim \mathcal B(n, p)$ 进行比较：在该假设下，计算
出现与 $k$ 相比\emph{至少同样极端}的结果的概率。若该概率
非常小（常用约定为低于 $5\%$），则拒绝该假设；
否则，观察结果与该假设相容。阈值是一种
选择，而非定理 --- 统计学只量化风险，
是否接受该风险则由使用者决定。
\end{method}

\section{练习}

\begin{exercise}[$\star$]\label{exo:g11:binom:1}
将帕斯卡三角形扩展到第 $7$ 行，并求
$\binom62$、$\binom{7}{3}$ 和 $\binom74$ 的值。
\end{exercise}

\begin{exercise}[$\star$]\label{exo:g11:binom:2}
将一颗均匀骰子掷 $4$ 次，$X$ 为掷出六点的次数。说明
$X \sim \mathcal B\left(4, \frac16\right)$，并计算 $\P(X = 0)$、
$\P(X = 1)$ 和 $\P(X \geq 2)$。
\end{exercise}

\begin{exercise}[$\star$]\label{exo:g11:binom:3}
下列各量中，哪些服从二项分布？说明理由。
\begin{enumerate}
  \item 掷一枚均匀硬币 $20$ 次，正面出现的次数。
  \item 从一副牌中抽取 $5$ 张牌，其中 A 牌的张数。
  \item 下周的雨天数，若每天下雨的概率为
        $0.3$，且各天相互独立。
\end{enumerate}
\end{exercise}

\begin{exercise}[$\star$]\label{exo:g11:binom:4}
设 $X \sim \mathcal B(50, 0.2)$，求 $\E(X)$、$\V(X)$ 和 $\sigma(X)$。
\end{exercise}

\begin{exercise}[$\star\star$]\label{exo:g11:binom:5}
一名射手每次射击击中目标的概率为 $0.7$，
各次射击相互独立。射击 $6$ 次，求恰好击中 $4$ 次的概率，
以及至少击中 $5$ 次的概率。
\end{exercise}

\begin{exercise}[$\star\star$]\label{exo:g11:binom:6}
一份判断题测验共 $8$ 题，某学生对每题均随机猜测答案。
求其及格（至少答对 $6$ 题）的概率。
\end{exercise}

\begin{exercise}[$\star\star$]\label{exo:g11:binom:7}
每个购买的麦片盒中装有小人偶 A 或小人偶 B 之一，
二者概率各为 $\frac12$，且相互独立。收藏者购买 $5$ 盒，
求收藏者至少得到每种小人偶各一个的
概率。（对立事件：全为 A 或全为 B。）
\end{exercise}

\begin{exercise}[$\star\star$]\label{exo:g11:binom:8}
一名篮球运动员每次罚球命中的概率为 $p$，
各次罚球相互独立。设 $X \sim \mathcal B(3, p)$ 为
三次罚球中命中的次数。将 $\P(X = 3)$ 和 $\P(X \geq 1)$ 表示为 $p$ 的函数，
并求 $p$ 为何值时，三次罚球全部命中的概率
等于 $\frac{27}{64}$。
\end{exercise}

\begin{exercise}[$\star\star$]\label{exo:g11:binom:9}
一枚均匀硬币至少需要掷多少次，才能使得
至少出现一次正面的概率超过 $0.99$？（先考虑对立事件，再逐个试验
$n$ 的取值。）
\end{exercise}

\begin{exercise}[$\star\star$]\label{exo:g11:binom:10}
利用帕斯卡法则（见\cref{prop:g11:binom:pascal}）及
$\binom n0 = \binom nn = 1$，证明帕斯卡三角形的每一行中各数之和
恰为 $2^n$：将等式两边都解释为计算树状图中的
全部路径条数。
\end{exercise}

\begin{exercise}[$\star\star\star$]\label{exo:g11:binom:11}
某政客声称其支持率为 $60\%$。在一个含 $10$ 人的随机样本中，
仅有 $3$ 人表示支持。
\begin{enumerate}
  \item 若该宣称成立，样本中表示支持的
        人数 $X$ 服从什么分布？求 $\P(X \leq 3)$。
  \item 利用\cref{met:g11:binom:decision}中的判定准则，取
        阈值为 $5\%$，该观察结果与该宣称是否相容？
\end{enumerate}
\end{exercise}

\section{问题：高尔顿板}

\begin{problem}[{周末问题 --- 小球、钉子与帕斯卡
三角形：钟形曲线如何形成，为何系列赛有利于较强的一方，
以及何时应当提出质疑}]
\label{pb:g11:binom:1}
让一千个小球穿过一排钉子组成的格阵，每次弹跳都相当于一次
左右等可能的抛硬币，下方的格子便会堆积成一个
平滑对称的钟形 --- 每次皆然。这台装置称为
高尔顿板，其数学恰是本章的
二项分布（见\cref{thm:g11:binom:binomial}）。本题
将构造三角形、运行高尔顿板、裁判一场
七场四胜制的系列赛，最后落在二项分布大显身手之处：
判断一个观察结果何时应使我们
对某一宣称产生怀疑。

\medskip
\noindent\textbf{第一部分 --- 三角形。}
\begin{enumerate}
  \item 作出帕斯卡三角形，直至第 $6$ 行
        （\cref{prop:g11:binom:pascal}）。用一句话叙述并解释
        对称性 $\binom nk = \binom{n}{n-k}$，
        （选出 $k$ 个对象，等同于\,\dots）。
  \item 在第 $4$ 行和第 $5$ 行上验证每一行之和均为 $2^n$，
        并加以证明：所有 $\binom nk$ 合起来
        计数的是什么？
  \item 用委员会论证法重新推导帕斯卡法则
        $\binom{n+1}{k} = \binom nk + \binom{n}{k-1}$，
        即先固定一位特定的人，
        再按此人的去留将委员会分类。
  \item 计算 $\binom73$ 两次：一次由三角形得出，
        另一次由阶乘公式得出。
  \item 验证阶梯恒等式
        $\binom22 + \binom32 + \binom42 + \binom52 =
        \binom63$，并通过
        自 $\binom63$ 向下逐次应用帕斯卡法则来解释。
\end{enumerate}

\medskip
\noindent\textbf{第二部分 --- 高尔顿板。}
一个小球穿过 $n$ 行钉子下落；每遇到一颗钉子，它便
以概率 $\frac12$ 向左或向右弹开，且各次弹跳相互独立。
按向右弹跳的次数，将各槽编号为 $0$ 至 $n$。
\begin{enumerate}[resume]
  \item 借助
        \cref{met:g11:binom:recognize}所列的检验清单，说明槽号为何
        服从二项分布
        $\mathcal B\!\left(n, \frac12\right)$。
  \item 对较小的板（$n = 4$），求出五个槽的
        概率。哪一个槽最为拥挤？
  \item 取 $n = 10$，共 $1\,024$ 个球，求中间槽、第 $7$ 号槽
        及各边缘槽中的\emph{期望}球数各为多少？
        并描述球堆的形状。
  \item 设 $X \sim \mathcal B\!\left(10, \frac12\right)$：
        求 $\E(X)$、$V(X)$ 及 $\sigma$
        （见\cref{prop:g11:binom:expectation}），再求
        期望落在中心 $2\sigma$ 之内的球的比例，
        （即槽 $2$ 至 $8$），并将其与
        由\cref{pb:g11:stat:1}给出的切比雪夫保证作比较。
  \item 倾斜的板以概率 $0.6$ 向右弹开：
        求 $\E$、$V$ 与 $\sigma$（$n = 10$），并描述
        球堆会发生怎样的变化。
  \item 用一两句话回答：板的设计中，是什么
        造就了钟形 --- 以及为何现实中的许多量
        （如身高、测量误差）同样以这种方式堆积？
        （两者背后的深刻定理是
        中心极限定理，它是
        大学分册概率课程的巅峰。）
\end{enumerate}

\medskip
\noindent\textbf{第三部分 --- 七场四胜。}
两支队伍进行一场系列赛：先赢得 $4$ 场者夺得冠军；
各场比赛相互独立。
\begin{enumerate}[resume]
  \item 两队实力相当（$p = \frac12$），求
        系列赛恰以横扫（恰好 $4$ 场）结束的概率。
  \item 求系列赛打满
        $7$ 场的概率（前 $6$ 场之后比分须为何？）。
  \item 补全系列赛长度的分布，
        即 $4$、$5$、$6$ 或 $7$ 场（两队实力相当），并
        求期望长度。哪一种长度
        最有可能出现？
  \item 现设一队每场获胜的概率为 $p = 0.6$，求该队
        夺得系列赛的概率（在 $4$、$5$、$6$
        或 $7$ 场内获胜，即每种情形中该队都在最后一场获胜，
        且此前已赢得其中 $3$ 场）。系列赛对
        每场优势有何影响？
  \item 将其与单场决赛（胜率 $60\,\%$）及三场两胜制
        （需计算其概率）相比较。说明系列赛长度
        对实力与运气的相对作用有何一般影响 --- 以及为何联赛
        偏好较长的决赛？
\end{enumerate}

\medskip
\noindent\textbf{第四部分 --- 何时应当提出质疑。}
\begin{enumerate}[resume]
  \item 一枚硬币抛掷 $100$ 次，出现 $62$ 次正面。对于一枚均匀硬币，
        求 $\E$、$\sigma$ 及该观察结果的
        标准分数（见\cref{pb:g11:stat:1}）。
        按 $2\sigma$ 约定，结论如何？
  \item 某供应商宣称其零件的次品率为 $2\,\%$。从一批
        $50$ 个零件中你发现 $3$ 个次品。求
        该宣称成立时的 $\P(X \geq 3)$
        （$X \sim \mathcal B(50, 0.02)$；逐一计算
        $\P(X = 0), \P(X = 1), \P(X = 2)$），在
        $5\,\%$ 阈值（见\cref{met:g11:binom:decision}，
        \cref{exo:g11:binom:11}）下，是否令人警觉？
  \item 彩票中的持久性问题：每张彩票中奖（某种奖项）的概率为
        $\frac{1}{1000}$。求
        用 $1\,000$ 张彩票时至少中奖一次的概率。答案
        （$\approx 63\,\%$，而非 $100\,\%$！）隐藏着一个著名常数：
        计算 $0.999^{1000}$，并牢记这一数值
        $0.368$，留待十二年级使用。
  \item 压轴题 --- 二项分布的小像：识别清单
        （$n$ 固定，各次独立，$p$ 恒定）；
        以帕斯卡三角形为其表格，以钟形为其形状；
        以 $np$ 与 $np(1-p)$ 为其罗盘；其两位继承者
        正等待着十二年级的登场 --- 光滑的钟形曲线与
        大数定律。每项各用一句话作答。
\end{enumerate}
\end{problem}
